\documentclass[11pt]{article}
\usepackage[margin=1in]{geometry}
\usepackage{amsmath,amssymb,amsthm}
\usepackage{hyperref}
\usepackage{enumitem}

\newtheorem{theorem}{Theorem}
\newtheorem{lemma}{Lemma}
\newtheorem{corollary}{Corollary}
\newtheorem{remark}{Remark}

\newcommand{\N}{\mathbb N}
\newcommand{\R}{\mathbb R}
\newcommand{\ol}{\overline}
\newcommand{\co}{\operatorname{co}}
\newcommand{\supp}{\operatorname{supp}}
\newcommand{\Ball}{B}

\title{A positive \texorpdfstring{$\ell_p$}{ell-p}-sum class for fully Mackey complete duals}
\author{Run \texttt{fa\_banach\_001}, agent \texttt{agent\_lane\_05}}
\date{June 14, 2026}

\begin{document}
\maketitle

\section*{Status}

This packet records a claimed partial positive answer to Problem 4.15 of
Guirao, Martinez-Cervantes and Rodriguez~\cite{GMR}.  The claim is not a full
solution of the problem; it proves the desired implication for a broad
nonseparable class of $\ell_p$-sums.

\section*{Source problem}

Guirao, Martinez-Cervantes and Rodriguez~\cite[Problem 4.15]{GMR} ask:

\begin{quote}
Is $X^*$ fully Mackey complete whenever $X$ is $w^*$-sequentially dense in
$X^{**}$?
\end{quote}

They prove the affirmative answer when $X$ is separable
\cite[Corollary 4.8]{GMR}, and when $X$ is Asplund
\cite[Corollary 4.9]{GMR}.  Their Remark 4.7 notes that a negative answer
would have to use a nonseparable, non-Asplund space without subspaces
isomorphic to $\ell_1$.

\section*{Proof Intuition}

The obstruction in the nonseparable case is that a norming closed subspace
$Y\subset X^{**}$ need not project cleanly onto the countable coordinates on
which a given $u\in X^{**}$ is supported.  The main observation is that, for
$\ell_p$-sums of separable spaces, one can enlarge the countable support of
$u$ to another countable set $A$ so that the actual intersection
$Y\cap X_A^{**}$ is still norming for $X_A^*$.  This is done recursively.
At each countable stage, the dual ball of the countable sub-sum is weak-star
compact metrizable, so a countable family of elements of $Y$ suffices to keep
the norming inequality on that stage; adding their countable supports prepares
the next stage.

Once this localization is available, the problem becomes separable.  The
countable sub-sum contains no $\ell_1$, so the source paper's separable
criterion says its dual is fully Mackey complete.  The resulting compact
$LK$ witness lives entirely on the countable coordinates and is therefore also
a valid witness in the original large sum.

\section*{Notation and source facts}

For a Banach space $Z$, write $Z$ as canonically embedded in $Z^{**}$.
For a dual pair $\langle Z^*,Z^{**}\rangle$, weak-star means
$\sigma(Z^{**},Z^*)$.

We use the characterization from~\cite[Theorem 3.5]{GMR}: a Banach space $W$
is fully Mackey complete if and only if it has property $(LK)$:
for every norming norm-closed subspace $Y\subset W^*$ and every
$u\in W^*\setminus Y$, there is an absolutely convex weak-star compact
$K\subset W^*$ such that
\[
K\subset Y\oplus [u],
\qquad
u\in \ol{K\cap Y}^{\,w^*}.
\]

We also use~\cite[Corollary 4.8]{GMR}: if $Z$ is separable, then $Z^*$ is
fully Mackey complete if and only if $Z$ contains no copy of $\ell_1$.

\section*{Main theorem}

\begin{theorem}\label{thm:main}
Let $1<p<\infty$, let $\Gamma$ be a set, and let
$(E_\gamma)_{\gamma\in\Gamma}$ be separable Banach spaces containing no
subspace isomorphic to $\ell_1$.  Put
\[
X=\left(\bigoplus_{\gamma\in\Gamma} E_\gamma\right)_{\ell_p}.
\]
Then $X$ is weak-star sequentially dense in $X^{**}$, and $X^*$ is fully
Mackey complete.
\end{theorem}

\begin{remark}
If $\Gamma$ is uncountable and at least one $E_\gamma$ is separable
non-Asplund without $\ell_1$, for example a James-tree type space, then this
gives a nonseparable and non-Asplund positive class for Problem 4.15.
\end{remark}

\section*{Preliminary lemmas}

For $A\subset\Gamma$, write
\[
X_A=\left(\bigoplus_{\gamma\in A}E_\gamma\right)_{\ell_p}.
\]
We identify $X_A$, $X_A^*$, and $X_A^{**}$ with the corresponding coordinate
subspaces of $X$, $X^*$, and $X^{**}$.  Let $P_A:X^{**}\to X_A^{**}$ be the
coordinate projection.  Every element of $X^{**}$ has countable support because
$1<p<\infty$.

\begin{lemma}\label{lem:countable-no-l1}
If $A\subset\Gamma$ is countable, then $X_A$ is separable and contains no copy
of $\ell_1$.
\end{lemma}

\begin{proof}
Separability follows from countability of $A$ and separability of each
$E_\gamma$.

To rule out $\ell_1$, use Rosenthal's $\ell_1$ theorem.  Let $(x_n)$ be a
bounded sequence in $X_A$.  Enumerate $A=\{\gamma_1,\gamma_2,\ldots\}$.  Since
each $E_{\gamma_j}$ contains no $\ell_1$, Rosenthal's theorem gives, by a
diagonal argument, a subsequence $(x_{n_k})$ such that for each $j$ the
coordinate sequence $(P_{\{\gamma_j\}}x_{n_k})_k$ is weakly Cauchy in
$E_{\gamma_j}$.

For any $x^*=(x_\gamma^*)_{\gamma\in A}\in X_A^*$, approximate $x^*$ in norm
by a finitely supported functional.  The finite-coordinate part is Cauchy on
$(x_{n_k})$, and the tail is uniformly small because $(x_{n_k})$ is bounded.
Thus $(x^*(x_{n_k}))_k$ is Cauchy for every $x^*\in X_A^*$, so $(x_{n_k})$ is
weakly Cauchy.  Rosenthal's theorem now implies that $X_A$ contains no
isomorphic copy of $\ell_1$.
\end{proof}

\begin{lemma}\label{lem:localization}
Let $Y\subset X^{**}$ be a norming norm-closed subspace.  For every countable
$A_0\subset\Gamma$ there is a countable $A\supset A_0$ such that
\[
Y_A:=Y\cap X_A^{**}
\]
is norming for $X_A^*$.
\end{lemma}

\begin{proof}
Since $Y$ is norming, choose $0<\delta\leq 1$ such that
\[
\sup\{|y(x^*)|:\ y\in Y,\ \|y\|\leq 1\}\geq \delta \|x^*\|
\qquad (x^*\in X^*).
\tag{1}
\]

We recursively build countable sets
\[
A_0\subset A_1\subset A_2\subset\cdots
\]
and countable sets $C_n\subset Y\cap \Ball_{X^{**}}$.
Assume $A_n$ is countable.  Then $X_{A_n}$ is separable, so
$(\Ball_{X_{A_n}^*},w^*)$ is compact metrizable.  For each rational
$r\in(0,1)$, put
\[
O_r=\{x^*\in \Ball_{X_{A_n}^*}:\ \|x^*\|>r\}.
\]
For every $x^*\in O_r$, (1) gives $y\in Y\cap\Ball_{X^{**}}$ with
$|P_{A_n}y(x^*)|=|y(x^*)|>\delta r$.  Hence the weak-star open sets
\[
\{z^*\in \Ball_{X_{A_n}^*}: |P_{A_n}y(z^*)|>\delta r/2\},
\qquad y\in Y\cap\Ball_{X^{**}},
\]
cover $O_r$.  Since $O_r$ is an open subspace of a compact metrizable space,
it is Lindelof.  Choose a countable subcover.  Taking the union over rational
$r\in(0,1)$ gives a countable set $C_n\subset Y\cap\Ball_{X^{**}}$ such that
\[
\sup_{y\in C_n}|P_{A_n}y(x^*)|\geq \frac{\delta}{2}\|x^*\|
\qquad (x^*\in X_{A_n}^*).
\tag{2}
\]
Now let $A_{n+1}$ be $A_n$ together with the supports of all elements of
$C_n$.  This is countable, because each element of $X^{**}$ has countable
support.

Set $A=\bigcup_n A_n$.  We claim that $Y_A=Y\cap X_A^{**}$ is norming for
$X_A^*$.  The case $x^*=0$ is trivial.  Let $x^*\in X_A^*$ be nonzero and
write $x_n^*$ for its restriction to $A_n$, extended by zero to $A$.  Then
$\|x^*-x_n^*\|\to0$.  Fix a small $\eta>0$ and choose $n$ so large that
\[
\|x^*-x_n^*\|<\eta\|x^*\|,
\qquad
\|x_n^*\|>(1-\eta)\|x^*\|.
\]
By (2), choose $y\in C_n$ which is within $\eta\|x^*\|$ of the supremum at
$x_n^*$, so
\[
|P_{A_n}y(x_n^*)|>
\frac{\delta}{2}\|x_n^*\|-\eta\|x^*\|.
\]
Since $y\in C_n\subset Y\cap X_A^{**}$ and $\|y\|\leq1$,
\[
|y(x^*)|
\geq |P_{A_n}y(x_n^*)|-|y(x^*-x_n^*)|
>
\left(\frac{\delta}{2}(1-\eta)-2\eta\right)\|x^*\|.
\]
Taking $\eta$ sufficiently small, we get for instance
\[
\sup\{|y(x^*)|:\ y\in Y_A,\ \|y\|\leq1\}\geq \frac{\delta}{4}\|x^*\|.
\]
Thus $Y_A$ is norming for $X_A^*$.
\end{proof}

\section*{Proof of Theorem \ref{thm:main}}

First we show that $X$ is weak-star sequentially dense in $X^{**}$.  Let
$u\in X^{**}$ and let $A=\supp(u)$, which is countable.  By
Lemma~\ref{lem:countable-no-l1}, $X_A$ is separable and contains no copy of
$\ell_1$.  The source paper's separable criterion~\cite[Corollary 4.8]{GMR}
implies that $(\Ball_{X_A^{**}},w^*)$ is Frechet-Urysohn.  Since $X_A$ is
weak-star dense in $X_A^{**}$ by Goldstine's theorem, a sequence from $X_A$
converges weak-star to $u$ after scaling to a suitable ball.  Therefore $X$ is
weak-star sequentially dense in $X^{**}$.

It remains to prove that $X^*$ is fully Mackey complete.  We use property
$(LK)$ from~\cite[Theorem 3.5]{GMR}.  Let $Y\subset X^{**}$ be a norming
norm-closed subspace and let $u\in X^{**}\setminus Y$.  Let
$A_0=\supp(u)$, and apply Lemma~\ref{lem:localization} to obtain a countable
$A\supset A_0$ such that $Y_A=Y\cap X_A^{**}$ is norming for $X_A^*$.

By Lemma~\ref{lem:countable-no-l1}, $X_A$ is separable and contains no
$\ell_1$.  Hence $X_A^*$ is fully Mackey complete by
\cite[Corollary 4.8]{GMR}.  Since $u\in X_A^{**}\setminus Y_A$, property
$(LK)$ for $X_A^*$ gives an absolutely convex weak-star compact set
$K\subset X_A^{**}$ such that
\[
K\subset Y_A\oplus [u],
\qquad
u\in \ol{K\cap Y_A}^{\,\sigma(X_A^{**},X_A^*)}.
\]
The zero-extension embedding $X_A^{**}\hookrightarrow X^{**}$ is weak-star to
weak-star continuous, and the weak-star topology induced from $X^{**}$ on the
coordinate subspace $X_A^{**}$ is exactly $\sigma(X_A^{**},X_A^*)$.  Therefore
the same set $K$, regarded as a subset of $X^{**}$, is absolutely convex and
weak-star compact.  Moreover
\[
K\subset Y\oplus [u],
\qquad
u\in \ol{K\cap Y}^{\,\sigma(X^{**},X^*)},
\]
because $Y_A\subset Y$.  This is precisely the $(LK)$ witness required for
$X^*$.  Thus $X^*$ is fully Mackey complete.

\section*{Verification and limitations}

The result is a partial positive answer to Problem 4.15, not a full solution.
It covers $\ell_p$-sums, $1<p<\infty$, of separable no-$\ell_1$ spaces.  The
proof uses the countable-coordinate support structure of $\ell_p$-sums in an
essential way.  It does not address arbitrary nonseparable, non-Asplund spaces
without $\ell_1$.

The main point for review is Lemma~\ref{lem:localization}.  The Lindelof
argument selects countably many elements of the original norming subspace at
each countable coordinate stage; adding their supports ensures that, in the
limit, those same elements lie in $Y\cap X_A^{**}$ and still norm the enlarged
countable sub-sum up to a fixed constant.

The bounded novelty search for this run found no later literature answer to
Problem 4.15 using the phrases ``fully Mackey complete'', ``fully Mazur'', and
the exact wording of the problem.  The local parsed corpus only showed the
source paper and one later bibliography citation for the exact terminology.

\section*{Human-review recommendation}

Treat this as a likely valid partial theorem.  If the localization lemma
survives review, this packet gives a new positive nonseparable class for
Problem 4.15, including examples outside the separable and Asplund cases
explicitly listed in the source paper.

\begin{thebibliography}{9}

\bibitem{GMR}
A. J. Guirao, G. Martinez-Cervantes, and J. Rodriguez,
\emph{Completeness in the Mackey topology by norming subspaces},
J. Math. Anal. Appl. \textbf{478} (2019), no. 2, 776--789;
arXiv:1812.10669.

\bibitem{Rosenthal}
H. P. Rosenthal,
\emph{A characterization of Banach spaces containing $\ell_1$},
Proc. Nat. Acad. Sci. U.S.A. \textbf{71} (1974), 2411--2413.

\end{thebibliography}

\end{document}
